% IEEE conference manuscript skeleton.
%
% Numbers come from % Generated by src/scripts/08_paper_numbers.py -- do not edit by hand.
% Re-run after changing parameters or re-running experiments.

\newcommand{\StemLength}{0.85}
\newcommand{\StemNodes}{16}
\newcommand{\RadiusBase}{5.0}
\newcommand{\RadiusTip}{2.0}
\newcommand{\YoungsModulus}{1.0}
\newcommand{\RigidityRatio}{23}
\newcommand{\ParticlesPerPlant}{25}
\newcommand{\ConstraintsPerPlant}{46}
\newcommand{\Substeps}{32}
\newcommand{\Iterations}{4}
\newcommand{\WindSpeed}{4.0}

\newcommand{\TipDeflectionMean}{0.136}
\newcommand{\TipDeflectionPeak}{0.287}
\newcommand{\TipDeflectionPercent}{34}
\newcommand{\TipHeight}{0.821}

\newcommand{\BenchMaxPlants}{16384}
\newcommand{\BenchMaxParticles}{409600}
\newcommand{\BenchWarpCPU}{2516.8}
\newcommand{\BenchRepeats}{3}
\newcommand{\BenchWarpCUDA}{116.04}
\newcommand{\Speedup}{21.7}
\newcommand{\GPUAvailable}{true}

\newcommand{\FieldValidationDone}{false}
\newcommand{\FieldN}{\textbf{??}}
\newcommand{\FieldMAE}{\textbf{??}}
\newcommand{\FieldSpecies}{\textbf{??}}
% WARNING: no field measurements yet. Run 06_field_validation.py.
 and % Generated by src/scripts/07_parameter_table.py -- do not edit by hand
\begin{table}[t]
\caption{Simulation parameters. Material constants are taken from published ranges for green herbaceous tissue; geometry describes a row crop at working height.}
\label{tab:parameters}
\centering
\begin{tabular}{llll}
\toprule
Symbol & Value & Unit & Basis \\
\midrule
g & 0, -9.81, 0 & m/s$^2$ & standard gravity \\
L & 0.85 & m & row crop at working height \\
n & 16 & -- & converged at this resolution; see VALIDATION.md \\
r\_base & 5.0 & mm & calliper-measurable; sets EI at the clamp \\
r\_tip & 2.0 & mm & taper gives EI $\propto$ r$^4$ falloff with height \\
E\_stem & 1.00 & GPa & green herbaceous stem, published range $\approx$0.3--5 GPa; values below $\approx$0.2 GPa cannot support the stem's own weight (see material sanity check) \\
$\rho$\_stem & 700 & kg/m$^3$ & fresh tissue, mostly water \\
Cd\_stem & 1.10 & -- & bluff cylinder in cross-flow \\
E\_leaf & 0.10 & GPa & lamina $\approx$10$\times$ more compliant than stem \\
L\_leaf & 0.140 & m & blade length \\
w\_leaf & 0.050 & m & sets drag area \\
Cd\_leaf & 1.40 & -- & flat plate normal to flow \\
U & 4.0 & m/s & moderate breeze at reference height \\
$\alpha$ & 0.20 & -- & power-law boundary layer over cropland \\
$\rho$\_air & 1.225 & kg/m$^3$ & ISA sea level \\
$\Delta$t & 0.01667 & s & 60 Hz frame \\
n\_sub & 32 & -- & converged; see convergence study \\
n\_iter & 4 & -- & converged; see convergence study \\
c & 1.20 & 1/s & $\zeta$$\approx$0.1 at ~1.5 Hz; rate, not per-frame factor \\
\bottomrule
\end{tabular}
\end{table}
, both
% generated from the code by src/scripts/08_paper_numbers.py and
% 07_parameter_table.py. Do not type a numeric result directly into this
% file -- if it is worth quoting, it is worth generating.
%
% Unfilled results render as bold ?? in the PDF. Search the compiled output
% for "??" before submitting.

\documentclass[conference]{IEEEtran}
\IEEEoverridecommandlockouts

\usepackage{cite}
\usepackage{amsmath,amssymb}
\usepackage{graphicx}
\usepackage{booktabs}
\usepackage{url}

% Generated by src/scripts/08_paper_numbers.py -- do not edit by hand.
% Re-run after changing parameters or re-running experiments.

\newcommand{\StemLength}{0.85}
\newcommand{\StemNodes}{16}
\newcommand{\RadiusBase}{5.0}
\newcommand{\RadiusTip}{2.0}
\newcommand{\YoungsModulus}{1.0}
\newcommand{\RigidityRatio}{23}
\newcommand{\ParticlesPerPlant}{25}
\newcommand{\ConstraintsPerPlant}{46}
\newcommand{\Substeps}{32}
\newcommand{\Iterations}{4}
\newcommand{\WindSpeed}{4.0}

\newcommand{\TipDeflectionMean}{0.136}
\newcommand{\TipDeflectionPeak}{0.287}
\newcommand{\TipDeflectionPercent}{34}
\newcommand{\TipHeight}{0.821}

\newcommand{\BenchMaxPlants}{16384}
\newcommand{\BenchMaxParticles}{409600}
\newcommand{\BenchWarpCPU}{2516.8}
\newcommand{\BenchRepeats}{3}
\newcommand{\BenchWarpCUDA}{116.04}
\newcommand{\Speedup}{21.7}
\newcommand{\GPUAvailable}{true}

\newcommand{\FieldValidationDone}{false}
\newcommand{\FieldN}{\textbf{??}}
\newcommand{\FieldMAE}{\textbf{??}}
\newcommand{\FieldSpecies}{\textbf{??}}
% WARNING: no field measurements yet. Run 06_field_validation.py.


\begin{document}

\title{Wind-Driven Biomechanics of Herbaceous Row Crops\\
for Agricultural Robot Simulation}

\author{\IEEEauthorblockN{Sudharshan S}
\IEEEauthorblockA{Karunya Institute of Technology and Sciences\\
Coimbatore, India\\
sudharshans@karunya.edu.in}
% ADVISOR: confirm co-authorship and ordering before submission.
}

\maketitle

\begin{abstract}
% Write last. Structure: agricultural robots are trained against rigid plants;
% wind-driven deformation is the load case that matters for perception and is
% absent from existing deformable-plant simulators; we model slender
% herbaceous stems as XPBD particle chains with bending compliance derived
% from a tapered EI profile; validated against a closed-form cantilever
% solution, an independent implementation, and measured deflection of a real
% plant.
%
% Do NOT claim GPU speed-up in the abstract unless \GPUAvailable is true.
\end{abstract}

\begin{IEEEkeywords}
agricultural robotics, simulation, deformable objects, position-based
dynamics, plant biomechanics
\end{IEEEkeywords}

\section{Introduction}
% Motivation: perception and manipulation both degrade when the simulated
% plant cannot move. Wind is the dominant everyday load on a field crop and
% continuously changes the plant's silhouette.
%
% State contributions precisely. Suggested, and defensible given the prior
% work in Section II:
%   1) a wind-loaded biomechanical model of slender herbaceous row crops,
%      with bending compliance derived from a tapered EI \propto r^4 profile
%      rather than tuned;
%   2) a solver-design study showing that chain-sequential parallelisation
%      is required for stem-like topologies, with the colour-parallel
%      alternative measurably \SI{7}{\times} too compliant at equal cost;
%   3) validation against a closed-form cantilever solution, an independent
%      second implementation, and a real plant.

\section{Related Work}
% See Research/RELATED_WORK.md for the verified summaries these paragraphs
% must be written from.
%
% MANDATORY: OrchardBench~\cite{munn2026orchardbench} and Gazebo
% Plants~\cite{deng2024gazeboplants} must both be cited and explicitly
% differentiated. They are the closest prior art and a reviewer will know
% them. Do not write "existing simulators treat plants as rigid" without
% immediately qualifying it -- cite Agri-Sim~\cite{agrisim2026} for the rigid
% -body claim, which is accurate for that platform.
%
% The honest separation: those works address woody trees and manipulator
% contact; this addresses slender herbaceous stems under aerodynamic load.

\section{Method}

\subsection{Plant Model}
% Particle chain; axial and bending constraints; leaves as compliant blades.
% The soil clamp: a pinned root fixes position but not tangent, so a buried
% anchor particle plus a root-hinged bending constraint is required.

The stem is discretised into \StemNodes{} particles over
\StemLength{}\,m, tapering from \RadiusBase{}\,mm radius at the soil to
\RadiusTip{}\,mm at the tip. Bending compliance follows from beam theory,
\begin{equation}
\alpha_\text{bend} = \frac{L^3}{4EI}, \qquad I = \frac{\pi r^4}{4},
\end{equation}
so that flexural rigidity falls with height as $r^4$; the base of the default
stem is \RigidityRatio{}$\times$ stiffer than the tip without any
height-dependent stiffness rule being imposed.

\subsection{Aerodynamic Loading}
% Drag on the relative velocity, and why a constant wind acceleration cannot
% reach equilibrium.
\begin{equation}
\mathbf{F} = \tfrac{1}{2}\rho C_d A \lVert \mathbf{u}-\mathbf{v} \rVert
             (\mathbf{u}-\mathbf{v})
\end{equation}

\subsection{Solver}
% XPBD; why compliance rather than iteration-count stiffness; the chain vs
% colour parallelisation result; substeps beat iterations.

% Generated by src/scripts/07_parameter_table.py -- do not edit by hand
\begin{table}[t]
\caption{Simulation parameters. Material constants are taken from published ranges for green herbaceous tissue; geometry describes a row crop at working height.}
\label{tab:parameters}
\centering
\begin{tabular}{llll}
\toprule
Symbol & Value & Unit & Basis \\
\midrule
g & 0, -9.81, 0 & m/s$^2$ & standard gravity \\
L & 0.85 & m & row crop at working height \\
n & 16 & -- & converged at this resolution; see VALIDATION.md \\
r\_base & 5.0 & mm & calliper-measurable; sets EI at the clamp \\
r\_tip & 2.0 & mm & taper gives EI $\propto$ r$^4$ falloff with height \\
E\_stem & 1.00 & GPa & green herbaceous stem, published range $\approx$0.3--5 GPa; values below $\approx$0.2 GPa cannot support the stem's own weight (see material sanity check) \\
$\rho$\_stem & 700 & kg/m$^3$ & fresh tissue, mostly water \\
Cd\_stem & 1.10 & -- & bluff cylinder in cross-flow \\
E\_leaf & 0.10 & GPa & lamina $\approx$10$\times$ more compliant than stem \\
L\_leaf & 0.140 & m & blade length \\
w\_leaf & 0.050 & m & sets drag area \\
Cd\_leaf & 1.40 & -- & flat plate normal to flow \\
U & 4.0 & m/s & moderate breeze at reference height \\
$\alpha$ & 0.20 & -- & power-law boundary layer over cropland \\
$\rho$\_air & 1.225 & kg/m$^3$ & ISA sea level \\
$\Delta$t & 0.01667 & s & 60 Hz frame \\
n\_sub & 32 & -- & converged; see convergence study \\
n\_iter & 4 & -- & converged; see convergence study \\
c & 1.20 & 1/s & $\zeta$$\approx$0.1 at ~1.5 Hz; rate, not per-frame factor \\
\bottomrule
\end{tabular}
\end{table}


\section{Validation}

\subsection{Closed-Form Beam Theory}
% Mesh refinement table: 9/17/33 nodes, converged, vs qL^4/(8EI).
% State the residual ~8% bias and its cause (elastic clamp, lumped mass).

\subsection{Independent Implementation}
% Warp kernels vs NumPy reference agree to ~5e-5 m over 300 frames.

\subsection{Measured Plant Deflection}
% FILL FROM results/field_validation.md once measurements exist.
% If \FieldValidationDone is false this section cannot be written, and the
% paper should not claim real-world validation.
% n = \FieldN{}, species \FieldSpecies{}, mean absolute error \FieldMAE{} deg.
%
% State the limitations explicitly and in the body text, not only in a
% limitations section: preliminary, small n, single species, and -- if the
% fan was uncalibrated -- that the airspeed scale was fitted rather than
% measured.

\section{Results}

\subsection{Wind Response}
The default plant reaches a mean tip deflection of
\TipDeflectionMean{}\,m, peaking at \TipDeflectionPeak{}\,m
(\TipDeflectionPercent{}\,\% of stem length) under a \WindSpeed{}\,m/s wind,
while the stem tip remains at \TipHeight{}\,m.

\subsection{Throughput}
% ONLY write this subsection if \GPUAvailable is true. Otherwise the honest
% move is to report CPU scaling and state that GPU evaluation is future work
% -- a GPU claim without GPU numbers, in a field where
% OrchardBench~\cite{munn2026orchardbench} reports measured GPU results, will
% not survive review.

\section{Limitations}
% Say these plainly; a reviewer finding them first is much worse.
% - material constants assumed from published ranges, not measured per plant
% - drag model, not a flow solve: no wake, no plant-to-plant sheltering
% - leaves are 1-D chains: no twist or furling
% - purely elastic: no plastic deformation or stem breakage
% - validation scope: preliminary, small n, single species

\section{Conclusion}

\section*{Acknowledgment}
% Advisor, institution, compute.

\bibliographystyle{IEEEtran}
\bibliography{refs}

\end{document}
